% ============================================================
% 一、引言 (Introduction)
% ============================================================
\section{引言}

\subsection{研究背景与动机}

在此处简要介绍应用背景：这个问题的实际来源（可以是图形学、视觉、机器学习、数据科学等领域的真实场景），为什么它是一个值得研究的数学问题。

\begin{itemize}
    \item 描述问题的应用场景和实际意义；
    \item 指出现有方法的局限性或尚未解决的理论问题；
    \item 明确本文要探讨的具体问题范围。
\end{itemize}

\subsection{相关工作}

简要回顾与该问题相关的已有工作，重点关注其中与本课程数学主题相关的部分。

\begin{itemize}
    \item 文献综述应聚焦于数学方法，而非仅是实验结果对比；
    \item 明确指出已有工作与本课程哪些章节相关；
    \item 引用应优先来自顶会/顶刊或 arXiv。
\end{itemize}

% 示例引用格式（删除或替换）：
% 相关方法最早由 \cite{Alpher02} 提出，后续工作 \cite{Alpher03,Alpher04} 从优化角度进行了改进。本课程讲义第 X 章详细讨论了该方法的核心思想。

\subsection{本文贡献}

用条目清晰列出本文的主要工作（3--4 点即可）：

\begin{enumerate}
    \item 对【具体问题】给出了严格的形式化定义，明确了变量、假设、目标与约束；
    \item 从【几何/物理/概率】视角给出了核心数学对象的直观解释；
    \item 推导了【具体算法或分析方法】，分析了其收敛性/复杂度/数值稳定性；
    \item 实现了核心算法并设计了可复现实验，验证了方法的有效性及边界条件。
\end{enumerate}

\subsection{论文结构}

简述后续各节的安排，与课程要求的能力主线（形式化→直观化→数学化→可计算）相呼应。
